\subsection{ניתוח התוצאות}

מן התוצאות עולה שזיהוי הרובוטים הוא רכיב חזק ויציב, ואילו זיהוי
הספרות הוא עדיין החלק הרגיש יותר. למרות זאת, המערכת כולה נשארת שימושית
מפני שההחלטה הסופית אינה נשענת על מקור מידע יחיד אלא על צבירת ראיות
לאורך זמן.

\subsection{מגבלות ומקרי קצה}

למרות ההישגים, למערכת עדיין יש כמה מגבלות ברורות:

\begin{itemize}
\item \textbf{תלות באיכות הצילום:} סרטונים מטושטשים, רחוקים מדי או
  כאלה שצולמו בזווית לא נוחה מקשים מאוד על קריאת הספרות ועל שמירת
  המעקב.

\item \textbf{חסימות בין רובוטים:} כאשר כמה רובוטים נצמדים או חוצים זה
  את זה, יש סיכון לאיבוד עקבות או להחלפת זהויות זמנית.

\item \textbf{היעדר \textenglish{Re-Identification} חזותי מלא:} המערכת
  הנוכחית אינה לומדת חתימה חזותית עמוקה של כל רובוט. לכן, לאחר חסימה
  ארוכה, היא מסתמכת בעיקר על הזיהוי מחדש של הספרות ועל לוגיקת ההצבעה.

\item \textbf{הערכת ביצועים חלקית:} חלק מן המדדים במערכת המלאה עדיין
  דורשים הערכה כמותית מסודרת יותר, למשל דיוק שיוך סופי על קבוצת בדיקה
  ייעודית או מדדי מעקב פורמליים.
\end{itemize}

מגבלות אלו אינן מבטלות את השימושיות של המערכת, אך הן מגדירות היטב את
גבולות היכולת הנוכחיים שלה.

דוגמאות וויזואליות לכישלונות של המודלים מצורפות בנספח \ref{app:success_failure} באיור \ref{fig:failure}.

\subsection{מיקום ביחס למצב הקיים}

הפרויקט אינו מתיימר להתחרות ישירות ב־\textenglish{state of the art}
האקדמי בזיהוי אובייקטים, במעקב או ב־\textenglish{OCR}. במקום זאת, הוא
מכוון לבעיה יישומית, ממוקדת ומוגדרת היטב: סקאוטינג אוטומטי של משחקי
\textenglish{FRC} מתוך סרטון וידאו מהיציע.

בהקשר זה, התרומה של הפרויקט היא בעיקר הנדסית ומערכתית:

\begin{itemize}
\item שילוב של שני מודלי זיהוי שונים בתוך צינור עיבוד אחד.
\item ניצול ידע דומייני חזק מאוד מן הלו"ז של המקצה.
\item המרה של פלטי ראייה ממוחשבת למסלולי תנועה ולייצוג שימושי לסקאוטינג.
\item חיבור בין צד שרת, ממשק משתמש ואלגוריתם ניתוח בפלטפורמה אחת.
\end{itemize}

לכן, הערך של המערכת נמדד פחות בשאלה האם היא שוברת שיאי דיוק כלליים,
ויותר בשאלה האם היא פותרת בפועל בעיה אמיתית בעולם ה־\textenglish{FRC}.

\subsection{הצעות לשיפור עתידי}

אם היו לי עוד כמה חודשי פיתוח, יש כמה כיוונים שהייתי שמח להמשיך בהם:

\begin{itemize}
\item \textbf{איסוף מערך נתונים גדול וריאלי יותר:} במיוחד עבור ספרות
  בתנאי צילום קשים, כדי לשפר את ההכללה של מודל הספרות.

\item \textbf{הוספת רכיב \textenglish{Re-Identification}:} מודל חזותי
  נוסף שיוכל לסייע בשמירת זהות גם כאשר הספרות אינן נראות לאורך זמן.

\item \textbf{כיול גאומטרי של המגרש:} התאמת נקודות התמונה לקואורדינטות
  אמיתיות על המגרש, כך שהמסלולים ימדדו ביחידות מרחק אמיתיות ולא רק
  במרחב הפיקסלים.

\item \textbf{פרוטוקול הערכה מסודר יותר:} יצירת סט בדיקה מסומן היטב עם
  אמת מידה מלאה לשיוך קבוצות ולעקיבה, כדי לאפשר השוואה כמותית חזקה יותר
  בין גרסאות שונות של המערכת.

\item \textbf{אופטימיזציה לביצועים:} צמצום זמן הריצה ושיפור היכולת
  לעבוד קרוב יותר לזמן אמת גם על חומרה חלשה יותר.
\end{itemize}

\subsection{שיקולים אתיים ושימוש אחראי}

מבחינה אתית, זהו פרויקט בעל סיכון נמוך יחסית, משום שהוא מנתח משחקי
רובוטיקה ציבוריים ולא מידע רפואי, פיננסי או אישי רגיש. עם זאת, עדיין יש
כמה שיקולים חשובים:

\begin{itemize}
\item \textbf{אמינות ההחלטות:} אם המערכת טועה בזיהוי קבוצה או במסלול,
  היא עלולה להוביל למסקנות סקאוטינג שגויות. לכן חשוב להתייחס אליה ככלי
  מסייע ולא כתחליף מוחלט לשיפוט אנושי.

\item \textbf{שקיפות:} המשתמש צריך להבין שהמערכת מבצעת שיוך הסתברותי,
  ושבמקרים קשים ייתכן חוסר ודאות או שגיאה.

\item \textbf{שימוש הוגן:} מערכת כזו צריכה לשמש ככלי לניתוח ושיפור
  ביצועים, ולא ככלי להסקת מסקנות נחרצות על קבוצה בודדת על סמך נתון חלקי
  או סרטון באיכות נמוכה.
\end{itemize}

\subsection{מסקנה סופית}

לסיכום, הפרויקט מראה שאפשר לבנות מערכת שימושית לסקאוטינג אוטומטי גם
בלי לפתור באופן מושלם כל תת־בעיה בנפרד. השילוב בין זיהוי, מעקב, זיהוי
ספרות וידע מוקדם מספק תוצר מעשי לניתוח משחקי \textenglish{FRC}.
